\chapter{Le code, fichier par fichier}
\label{ann:code}

\begin{chapeau}
Cette annexe décrit chaque fichier du programme : son rôle, ses entrées, ses sorties et ses fonctions principales. Le
nombre de lignes est donné pour se faire une idée de la taille. Au total, le projet compte environ 10\,500 lignes de
Python et 6\,900 lignes de JavaScript et de HTML, dont environ 1\,900 lignes de tests et de campagnes.
\end{chapeau}

\newcommand{\fiche}[2]{\par\bigskip\noindent{\sffamily\bfseries\color{titre}\fichier{#1}}\hfill{\small\sffamily\color{gris}#2}\par\nopagebreak\smallskip}

\section{Arborescence}

\begin{lstlisting}[language={}, basicstyle=\ttfamily\footnotesize]
CT-CORA3D/
  CT-CORA3D.app            lanceur macOS de développement
  Lancer sur Windows.bat   lanceur Windows de développement
  LISEZ-MOI.md             mode d'emploi détaillé
  application/             serveur, fenêtre, interface, fabrication des applications
  code_source/             étapes 1 à 7 (Python), tests, cora/ (planification, JS)
  documents/               cahier des charges, rapport, documentation, ce dossier
  resultats/               un dossier par série, os fictif, rapports des campagnes
\end{lstlisting}

\section{Le programme de calcul : \texorpdfstring{\fichier{code_source/}}{code\_source/}}

\fiche{parametres.py}{198 lignes}
\textbf{Rôle} : tous les réglages du programme, au même endroit, chacun avec son origine (document de cadrage,
littérature ou mesure sur les scanners du projet). Aucune autre valeur « magique » ne doit apparaître dans les scripts.
Il définit aussi les dossiers : le dossier du projet, celui des scanners (\fichier{../CT-scan/} s'il existe) et celui
des résultats. Toutes les valeurs sont listées à l'\cref{ann:parametres}.

\fiche{etape1_lecture_dicom.py}{251 lignes}
\textbf{Rôle} : inventaire des séries d'un dossier et construction du volume (\cref{chap:lecture}).
\textbf{Entrée} : un dossier DICOM. \textbf{Sorties} : objets \code{Serie} et \code{Volume} (en mémoire).
\textbf{Fonctions} : \code{inventaire()} (fouille et regroupement par série), \code{charger()} (tri le long de la
normale, contrôles, rééchantillonnage, conversion en HU, matrice $\mt{A}$), \code{lire_pixels()},
\code{trouver()} (série par son nom) ; la classe \code{Volume} convertit indices et millimètres
(\code{vers_lps()}, \code{vers_indices()}).

\fiche{etape2_segmentation.py}{374 lignes}
\textbf{Rôle} : séparer les os (\cref{chap:segmentation}). \textbf{Entrée} : le \code{Volume}. \textbf{Sortie} :
un objet \code{Segmentation} (étiquettes, image lissée, géométrie de la zone recadrée).
\textbf{Fonctions} : \code{recadrage()}, \code{lisser()}, \code{graines_automatiques()} (seuil qui monte),
\code{croissance()} (ligne de partage des eaux), \code{contacts()}, \code{fusionner_surdecoupages()} (critère $c^*$),
\code{nettoyer()}, \code{segmenter()} (enchaîne le tout) ; outils de l'étape 3 : \code{recroissance()},
\code{raffiner()}, \code{boule()} ; \code{depuis_fichier()} relit une segmentation enregistrée ;
\code{boite_englobante()} calcule une boîte englobante 50 fois plus vite que la méthode classique.

\fiche{etape3_identification.py}{644 lignes}
\textbf{Rôle} : nommer les os et corriger la segmentation (\cref{chap:identification}). \textbf{Entrées} : la
\code{Segmentation}, les points de l'utilisateur. \textbf{Sorties} : noms, raisons, repère de la tête du fémur.
\textbf{Fonctions} : \code{decrire()} (descripteurs, ACP), \code{identifier()} (les 13 règles), \code{tete_femorale()}
et \code{ajuster_sphere()} (chapeau haut-de-forme et RANSAC), \code{separer_fibula()}, \code{separer_par_points()},
\code{nettoyer_morceaux()}, \code{nommer_cotes()}.

\fiche{etape4_maillage.py}{396 lignes}
\textbf{Rôle} : un STL étanche par os et la segmentation au format de 3D Slicer (\cref{chap:maillage}).
\textbf{Sorties} : \fichier{<os>.stl}, \fichier{segmentation.seg.nrrd}, les mesures de chaque maillage.
\textbf{Fonctions} : \code{exports()} (un fichier par os, plus le groupe « tarse »), \code{remplir()},
\code{mailler()}, \code{affiner_demi_hauteur()}, \code{controle()}, \code{bords_touches()}, \code{ecrire_stl()},
\code{ecrire_seg_nrrd()}, \code{lire_seg_nrrd()}, \code{exporter()}.

\fiche{etape5_controle.py}{267 lignes}
\textbf{Rôle} : sensibilité au seuil, aperçus et rapports (\cref{sec:sensibilite}). \textbf{Sorties} :
\fichier{rapport.md}, \fichier{rapport.json}, \fichier{apercu_3d.png}, \fichier{apercu_coupes.png},
\fichier{reperes_auto.mrk.json}. \textbf{Fonctions} : \code{sensibilite_seuil()}, \code{apercu_3d()},
\code{apercu_coupes()}, \code{markups()}, \code{rapport_md()}, \code{controler()}.

\fiche{etape6_reperes.py}{1236 lignes}
\textbf{Rôle} : repère, côté, complétude, repères anatomiques, angles et torsion du fémur et du tibia
(\cref{chap:reperes}). \textbf{Entrées} : les STL du dossier de résultats (et la fibula, la tête du fémur).
\textbf{Sorties} : \fichier{cora_<os>.json}, \fichier{cora_<os>.bin}, \fichier{cora_<os>.png}, \fichier{cora_cas.json}.
\textbf{Fonctions} : classe \code{Os} (repère courant, sections), \code{orienter()}, \code{cotes_femur()},
\code{cotes_tibia()}, \code{extremites_coupees()}, \code{analyser()} (les trois passes), \code{reperes_femur()},
\code{reperes_tibia()}, \code{tete_femorale()}, \code{centre_col()}, \code{petit_trochanter()},
\code{enveloppe_tangente()}, \code{sphere_robuste()}, \code{region_connexe()}, \code{angles_articulaires()},
\code{torsion()}, \code{calibration()}, \code{normales()}, \code{traiter_os()}, \code{choisir_cible()}. Ligne de
commande : \code{python etape6_reperes.py [dossiers] [--os tibia]}.

\fiche{etape7_guide.py}{2032 lignes}
\textbf{Rôle} : guide de coupe (deux moitiés reliées par des ponts sécables) et guide de réduction, avec leurs
contrôles (\cref{chap:guides}). \textbf{Entrées} : \fichier{prescription.json}, le maillage de l'os planifié, les os voisins et
le cartilage. \textbf{Sorties} : \fichier{guide_coupe.stl}, \fichier{guide_reduction.stl}, \fichier{guide_coupe_scanner.stl}, \fichier{os_corrige.stl},
\fichier{guide.json}, \fichier{guide.png}, les plans 2D, des \fichier{.bin} pour l'affichage.
\textbf{Fonctions} : \code{generer()}, \code{verifier_abord()}, \code{parametres_du_cas} (dimensions choisies pour un
cas), \code{Grille}, \code{distance_signee()}, \code{DistanceSurface}, \code{dist2_point_triangle()}, les champs
\code{champ_boite_arrondie()}, \code{champ_cylindre()}, \code{champ_capsule()}, \code{champ_secteur()}, \code{smin()},
\code{carte_hauteur()}, \code{depouille()}, \code{ombre()}, \code{EntreOs}, \code{sans_parties_fines()},
\code{isosurface()}, \code{guide_coupe()}, \code{transformation_reelle()}, \code{os_corrige()},
\code{guide_reduction()}, \code{controles_coupe()}, \code{controles_voisins()}, \code{controles_reduction()},
\code{protocole()},
\code{dimensions_guide()}, \code{assurer_jeu()}, \code{exporter()}.

\fiche{dessin_guide.py}{1230 lignes}
\textbf{Rôle} : les plans 2D du guide, générique (cotes en symboles) et coté (ce cas), en SVG, PDF et PNG, présentés
comme des plans de pièce industrielle (normes ISO). \textbf{Méthode} : projections orthogonales en projection
européenne, parties cachées par un tampon de profondeur en numpy, arêtes visibles en vecteurs ; coupe exacte
(manifold3d) hachurée ; cotes ISO, repères des pièces, nomenclature et cartouche avec le logo de l'ECAM
(\fichier{logo_ecam.png}) ; matplotlib, sans écran. \textbf{Fonctions} : \code{dessiner()}, \code{plan_generique()},
\code{plan_cote()}, \code{_rendu()} (tampon de profondeur), \code{_cote()}, \code{_hachures()},
\code{_cartouche()}, \code{_nomenclature()}.

\fiche{lancer_tout.py}{113 lignes}
\textbf{Rôle} : enchaîner les étapes 1 à 6 en ligne de commande, pour une série ou toutes. \textbf{Fonction} :
\code{traiter()}, appelée aussi par le serveur. Options : \code{--liste}, \code{--serie}, \code{--dossier},
\code{--sortie}, \code{--seuil}, \code{--ras}, \code{--sans-controle}, \code{--sans-images}.

\fiche{vtk_leger.py}{22 lignes}
\textbf{Rôle} : n'importer que les classes de VTK utilisées (0,5~s au lieu de 2,2~s au démarrage). Usage :
\code{import vtk_leger as vtk}.

\fiche{calibration_reperes.json}{fichier de données}
\textbf{Rôle} : cache des écarts de définition des angles (\cref{sec:calibration}), recalculé si
\fichier{etape6_reperes.py} change.

\section{Les tests : \texorpdfstring{\fichier{code_source/tests/}}{code\_source/tests/}}

\fiche{test_chaine.py}{172 lignes} Faux scanner écrit en DICOM, de géométrie connue ; contrôle de bout en bout des
étapes 1 à 5 (14 contrôles).

\fiche{test_reperes.py}{149 lignes} Os fictif dans des poses au hasard (repères exacts connus) et os du lévrier
(40 contrôles).

\fiche{test_guide.py}{378 lignes} Guides sur un os coudé synthétique (prescriptions écrites à la main) et sur les os du
lévrier ; option \code{--rapide}.

\fiche{modele_cora.py}{33 lignes} Lit l'os fictif (\fichier{cora/build/bones_pack.*}) depuis Python, pour les tests.

\fiche{campagne_export.mjs}{111 lignes} Génère les cas d'une campagne de guides (planification par le cœur
JavaScript, os, voisins et cartilage déformés).

\fiche{campagne_guides.py}{519 lignes} Construit les guides des cas en parallèle et les contrôle « avec l'œil du
chirurgien » ; écrit \fichier{rapport_campagne.md} et \fichier{resultats_campagne.jsonl}. Options \code{--processus},
\code{--filtre}, \code{--reprendre}, \code{--rapport}.

\fiche{campagne_strategie.mjs}{209 lignes} Campagne de la stratégie (préréglages et cas au hasard, règles R1 à R9) ;
option \code{--exporter} pour préparer une campagne de guides.

\section{La planification : \texorpdfstring{\fichier{code_source/cora/}}{code\_source/cora/}}

\fiche{app/cora_simple.js}{1727 lignes}
\textbf{Rôle} : le cœur de calcul de la planification, sans graphisme, utilisable sous Node (\cref{chap:planification}).
\textbf{Contenu} : vecteurs 2D et 3D, rotations (Rodrigues) et transformations rigides, réglages, préréglages
cliniques, paramètres du guide générique, valeurs publiées, \code{modeleReel()} (os segmenté au format du modèle),
\code{makeDeformity()}, \code{reference()}, \code{axesNormales()}, \code{coraPerpCommune()}, \code{planifier()},
\code{choisirAbord()}, \code{contraintesCoupe()}, \code{recommanderCoin()}, \code{placerBroches()},
\code{prescription()}, \code{mesuresTorsion()}, \code{planTorsion()}, \code{clipMesh()} (découpe de maillages),
\code{stlBinary()}.

\fiche{app/modules/10_base.js}{154 lignes} Utilitaires de l'interface (DOM, formatage), décodage du modèle osseux, et
la liste des références bibliographiques avec leurs DOI.

\fiche{app/modules/20_scene.js}{200 lignes} La scène 3D : rendu, caméra orthographique (comme une radiographie),
matériaux, et outils de dessin des tracés (lignes, points, arcs d'angle, plans, étiquettes).

\fiche{app/modules/30_os.js}{274 lignes} Les os dans la scène : os déformé, os sain en fantôme, fragments de
l'ostéotomie, broches, os tordu de l'étape « torsion ».

\fiche{app/modules/40_etapes.js}{1137 lignes} Les 13 étapes : titre, vue, légende et texte explicatif de chacune, avec
leurs chiffres et leurs sources.

\fiche{app/modules/45_export.js}{134 lignes} L'export pour le guide : archive ZIP (prescription, STL de l'os et des
voisins, carte du cartilage, mode d'emploi), écrite sans bibliothèque.

\fiche{app/modules/50_main.js}{1044 lignes} Le programme principal de l'onglet : état, panneaux, curseurs, sources (os
segmenté ou fictif), caméra, animation de la correction, construction des guides par le serveur, rendu à la demande
(\code{demanderImage()}).

\fiche{app/index.src.html}{401 lignes} Le gabarit de la page : structure et styles.

\fiche{build_app.py}{97 lignes} Assemble le gabarit, le cœur, les modules, Three.js, les polices et l'os fictif en un
seul fichier, \fichier{application/interface/cora/index.html}. À relancer après toute modification du JavaScript.

\fiche{generate_bones.py}{1149 lignes} Fabrique l'os fictif (fémur, tibia, fibula, rotule, fabelles d'un chien de
30~kg) par fonctions de distance et marching cubes, avec ses repères exacts et sa carte de cartilage ; écrit
\fichier{build/bones_pack.json} et \fichier{bones_pack.bin.gz}.

\fiche{tests/test_cas.mjs}{163 lignes} Planification sur l'os fictif : 12 cas, une grille d'amplitudes et de niveaux,
la torsion ; vérifie $\theta$, $F$, $S$, la formule de l'angle vrai, le CORA, la correction et les broches.

\fiche{tests/test_reel.mjs}{138 lignes} Planification sur les os segmentés, tels quels puis avec six déformités de test.

\section{L'application : \texorpdfstring{\fichier{application/}}{application/}}

\fiche{serveur.py}{1027 lignes} Le serveur local et les actions (inventaire, ouverture, segmentation, corrections,
guides) ; sert l'interface ; lance la fenêtre (\cref{sec:serveur}). Options : \code{--port}, \code{--sans-fenetre},
\code{--dossier}.

\fiche{fenetre.py}{288 lignes} La fenêtre native (pywebview) ; \code{--verifier} et \code{--installer} pour les
lanceurs.

\fiche{ct_cora3d_app.py}{80 lignes} Le point d'entrée de l'application autonome : sans option, la fenêtre ; avec
\code{--serveur}, le processus de calcul.

\fiche{interface/index.html}{130 lignes} La page à deux onglets (thème clair ou sombre, bouton Quitter).

\fiche{interface/ct/}{app.js 828, index.html 205 lignes} L'onglet \ui{Segmentation CT} : liste des séries, coupes,
os en 3D, corrections.

\fiche{interface/cora/index.html}{fichier produit} L'onglet \ui{Planification CORA}, assemblé par
\fichier{build_app.py}.

\fiche{lancer_mac.sh}{script} Choisit l'environnement Python (ou le crée au premier lancement) et lance le serveur ;
appelé par \fichier{CT-CORA3D.app}.

\fiche{lanceurs/}{scripts} \fichier{creer_app_autonome.sh} (application autonome macOS), \fichier{creer_app_mac.sh}
(petit lanceur \fichier{.app}), \fichier{creer_app_windows.bat} (application autonome Windows),
\fichier{ct_cora3d.spec} (recette PyInstaller commune), \fichier{lanceur_mac.applescript}.

\fiche{icone/creer_icone.py}{68 lignes} Dessine l'icône (coupe de scanner, fémur, axes et CORA) aux formats PNG, ICNS
et ICO.

\fiche{requirements.txt}{liste} Les bibliothèques Python : numpy, scipy, pydicom, SimpleITK, vtk, trimesh, pyvista,
matplotlib, pillow, manifold3d.

\section{Les documents : \texorpdfstring{\fichier{documents/}}{documents/}}

Le cahier des charges, le rapport de projet et la documentation technique (PDF et Word), leurs sources Markdown et
leurs figures, et ce dossier complet (\fichier{documents/dossier_complet/} : sources LaTeX, scripts des figures et
schémas draw.io modifiables).
